\documentclass[10pt,twocolumn]{article}

%% ── Geometry: balanced two-column margins ───────────────────────────────────
\usepackage[top=1in, bottom=1in, left=0.75in, right=0.75in,
            columnsep=0.25in]{geometry}

%% ── Core packages ───────────────────────────────────────────────────────────
\usepackage{amsmath,amssymb,amsthm}
\usepackage{booktabs}
\usepackage{array}
\usepackage{hyperref}
\usepackage{xcolor}
\usepackage{enumitem}
\usepackage{microtype}
\usepackage[T1]{fontenc}
\usepackage{lmodern}
\usepackage{titlesec}
\usepackage{caption}
\usepackage{multirow}
\usepackage{tabularx}
\usepackage{dblfloatfix} % allows table* on same page as text in twocolumn
\usepackage{flushend}   % balances columns on last page in twocolumn
\usepackage{listings}   % for code blocks (no overflow unlike verbatim)

%% ── parskip REMOVED — it breaks twocolumn column-height balance ─────────────

%% ── Listings style for code ─────────────────────────────────────────────────
\lstset{
  basicstyle=\ttfamily\footnotesize,
  breaklines=true,
  breakatwhitespace=false,
  columns=flexible,
  keepspaces=true,
  frame=none,
  xleftmargin=0pt,
  xrightmargin=0pt
}

%% ── Theorem environments ────────────────────────────────────────────────────
\newtheorem{theorem}{Theorem}[section]
\newtheorem{proposition}[theorem]{Proposition}
\newtheorem{corollary}[theorem]{Corollary}
\newtheorem{definition}[theorem]{Definition}
\newtheorem{remark}[theorem]{Remark}
\newtheorem{conjecture}[theorem]{Conjecture}

%% ── Custom commands ─────────────────────────────────────────────────────────
\newcommand{\cl}{\mathrm{cl}}
\newcommand{\SA}{\mathrm{SA}}
\newcommand{\Om}{\Omega}
\newcommand{\B}{\mathbf{B}}
\newcommand{\N}{\mathbb{N}}
\newcommand{\Z}{\mathbb{Z}}

\hypersetup{
  colorlinks=true,
  linkcolor=blue!60!black,
  citecolor=blue!60!black,
  urlcolor=blue!60!black,
}
% Allow URL line breaks everywhere
\def\UrlBreaks{\do\/\do\-\do\_\do\.\do\=\do\&\do\?}
\urlstyle{same}

%% ── Short alias for the long GitHub URL ─────────────────────────────────────
\newcommand{\repourl}{%
  \href{https://github.com/Prime-Spectrum-Team/Numbers-are-Information-.git}%
       {\small\texttt{github.com/Prime-Spectrum-Team/}\\\texttt{Numbers-are-Information-}}%
}

%% ── Title / Author / Date ───────────────────────────────────────────────────
\title{\textbf{Arithmetic Classification and Multiplicative Monoid
Structure of Natural Numbers: A Seven-Class Partition, Prime Coordinate
Representation, and Deterministic Addition Laws}}

\author{Prime-Spectrum-Team}

\date{Version 1.1 --- April 10, 2026}

%% ════════════════════════════════════════════════════════════════════════════
\begin{document}

\maketitle

\begin{abstract}
The natural numbers are conventionally treated as ordered magnitudes.
This paper develops a complementary framework: every natural number $n$
encodes intrinsic structural information through its prime
factorization---information that is independent of the quantity $n$
represents and computable in deterministic, finite time.

We make five contributions, three of which systematize classical results
into a unified framework and two of which are, to our knowledge, novel.

\emph{Systematization.}
\textbf{(C1)} We define a seven-class arithmetic partition of $\N$ based
on two invariants---the total prime factor count $\Om(n)$ and the residue
class of primes modulo 6---and prove it forms a finite commutative monoid
$(S, \otimes)$ under induced multiplication, with a complete 49-entry
Cayley table.
\textbf{(C2)} We characterize the power dynamics $\cl(n^k)$ and prove the
Universality Theorem: every $n > 1$ reaches the absorbing class $S_6$ in
at most 3 steps.
\textbf{(C5)} We define the resonance relation $R(i,j) = \Om(\gcd(i,j))$
and prove that $\approx 61\%$ of integer pairs are non-resonant (coprime),
matching the classical value $6/\pi^2$.

\emph{Novel contributions.}
\textbf{(C3)} We identify exactly 7 of the 28 unordered class pairs that
yield a deterministic addition law: six are trivially forced by the
singleton classes ($S_0=\{1\}$, $S_3=\{2\}$, $S_4=\{3\}$), and one---
Solar $+$ Lunar $\to S_6$---is the unique \emph{non-trivially}
deterministic law (Theorem~\ref{thm:solar-lunar}, verified on 90{,}000 pairs).
\textbf{(C4)} We formalize the SpectralAddress---a $K$-dimensional vector
of $p$-adic valuations, lossless for $\B$-smooth integers---and prove four
$O(K)$ isomorphisms mapping multiplication, GCD, LCM, and coprimality
testing to elementary vector operations.

All results are computationally verified in Python (100{,}000 pairs per
foundational axiom), C++20 ($\geq$475\,M Cayley lookups/s on Intel Core
i5; 559\,M verified on AMD Ryzen~9~3900), and a Verilog RTL
module intended for synthesis, validated via Python behavioral emulation
(14/14 test vectors; on-device synthesis and timing simulation pending).

\smallskip
\noindent\textbf{Keywords:} prime factorization, arithmetic
classification, multiplicative monoid, $p$-adic valuation,
SpectralAddress, resonance, deterministic addition laws, FPGA number
theory

\smallskip
\noindent\textbf{Repository:} \repourl
\end{abstract}

%% ────────────────────────────────────────────────────────────────────────────
\section{Introduction}
%% ────────────────────────────────────────────────────────────────────────────

\subsection{The Standard View and Its Limitation}

In computing, integers serve as universal indices: array positions, memory
addresses, and sequence positions are identified by integers whose sole
operationally relevant property is their ordinal position.

This view is not wrong, but it is incomplete.

The prime factorization of any integer $n > 1$ encodes structural
information that is:
\begin{itemize}[noitemsep, topsep=2pt]
  \item \textbf{Intrinsic}---present in the number itself;
  \item \textbf{Deterministic}---computable by a fixed procedure;
  \item \textbf{Compositional}---the structure of a product is determined
        by the structures of the factors.
\end{itemize}

Consider $n = 60 = 2^2 \times 3 \times 5$.  This factorization assigns it
depth $\Om(60) = 4$ and places it at the intersection of three prime axes.
These are not properties \emph{about} 60; they are properties 60
\emph{is}.

\subsection{Complementary Viewpoint}

The standard view treats $n$ as an ordinal. The complementary view
developed here treats $n$ as a coordinate in prime space: $60$ has
SpectralAddress $[2,1,1,0,0,0,0]$, placing it at the intersection of
prime axes 2, 3, and 5. Both views are consistent; they are adapted to
different purposes. This prime-arithmetic view has a concrete operational
consequence: every arithmetic operation on $\B$-smooth integers reduces to
a corresponding linear operation on SpectralAddress vectors---all in $O(K)$
time.

\subsection{Contributions}

\begin{itemize}[noitemsep, topsep=2pt]
\item[\textbf{C1}] Seven-class partition and multiplicative monoid
      (Theorems~\ref{thm:welldef}--\ref{thm:monoid}). \emph{Systematization.}
\item[\textbf{C2}] Power dynamics and universality (Theorem~\ref{thm:universality}). \emph{Systematization.}
\item[\textbf{C3}] Exactly 7 deterministic addition laws; 1 non-trivially
      deterministic (Theorem~\ref{thm:solar-lunar}). \emph{Novel.}
\item[\textbf{C4}] SpectralAddress and four $O(K)$ isomorphisms
      (Propositions~\ref{prop:mult-isom}--\ref{prop:coprime}). \emph{Novel.}
\item[\textbf{C5}] Resonance and $6/\pi^2$ sparsity (Theorem~\ref{thm:coprime-density}). \emph{Systematization.}
\end{itemize}

\subsection{Related Work}

The individual components---$\Om$-function~\cite{HardyWright1979,Tenenbaum2015},
prime factor count distributions~\cite{Sathe1953,Selberg1954},
$p$-adic valuations~\cite{Gouvea1997}---are classical.
A detailed positioning relative to prior work is given in Section~\ref{sec:relwork}.

%% ────────────────────────────────────────────────────────────────────────────
\section{Preliminaries}
%% ────────────────────────────────────────────────────────────────────────────

\subsection{The $\Omega$-Function}

\begin{definition}\label{def:omega}
For $n \in \N$ with canonical factorization $n = \prod_i p_i^{a_i}$,
define $\Om(n) = \sum_i a_i$ for $n > 1$ and $\Om(1) = 0$.
\end{definition}

\begin{proposition}[Complete additivity of $\Om$]\label{prop:omega-add}
For all $a, b \in \N$: $\Om(ab) = \Om(a) + \Om(b)$.
\end{proposition}
\begin{proof}
From unique prime factorization:
$\Om(ab) = \sum(\alpha_i + \beta_i) = \Om(a) + \Om(b)$. \qed
\end{proof}

\emph{This single axiom generates all 49 multiplication laws of the
monoid.} The Erd\H{o}s--Kac theorem~\cite{ErdosKac1940} establishes
that $\Om(n)$ is asymptotically normally distributed with mean and
variance $\ln \ln n$, confirming that $S_6$ dominates at large $N$.

\smallskip
\noindent\textbf{Verification (E1.1).} 100{,}000 random pairs from
$[1, 10^6]$: success rate 1.0 (95\% CI: [0.99996, 1.0]).

\subsection{Prime Residues Modulo 6}

\begin{proposition}\label{prop:prime-mod6}
Every prime $p > 3$ satisfies $p \equiv 1 \pmod{6}$ or
$p \equiv 5 \pmod{6}$.
\end{proposition}
\begin{proof}
Among $\{0,1,2,3,4,5\}$ mod~6: residues $\{0,2,4\}$ imply divisibility
by 2; residue 3 implies divisibility by 3. For prime $p > 3$, only
residues 1 and 5 remain. \qed
\end{proof}

By Dirichlet's theorem~\cite{Dirichlet1837}, these two residue classes
contain asymptotically equal numbers of primes. \emph{Verified on all
664{,}579 primes below $10^7$: ratio $= 0.992$.}

\subsection{The $p$-adic Valuation}

\begin{definition}\label{def:p-adic}
For prime $p$ and $n \in \N$, the \emph{$p$-adic valuation}
$\nu_p(n) = \max\{a \geq 0 : p^a \mid n\}$.
\end{definition}

Key property: $\nu_p(ab) = \nu_p(a) + \nu_p(b)$.

%% ────────────────────────────────────────────────────────────────────────────
\section{The Seven-Class Partition and Multiplicative Monoid}
%% ────────────────────────────────────────────────────────────────────────────

\subsection{Definition and Interpretation}

\begin{definition}[PrimeSpec classification]\label{def:primespec}
Define $\cl : \N \to S = \{S_0, S_1, S_2, S_3, S_4, S_5, S_6\}$:
\[
\cl(n) = \begin{cases}
S_0 & n = 1 \;\text{(Unit)} \\
S_3 & n = 2 \;\text{(Dyadic)} \\
S_4 & n = 3 \;\text{(Triadic)} \\
S_1 & \Om(n)=1,\; n \equiv 1\!\pmod{6} \;\text{(Solar)} \\
S_2 & \Om(n)=1,\; n \equiv 5\!\pmod{6} \;\text{(Lunar)} \\
S_5 & \Om(n) = 2 \;\text{(Semiprime)} \\
S_6 & \Om(n) \geq 3 \;\text{(Multiprime, absorbing)}
\end{cases}
\]
\end{definition}

\begin{table}[htbp]
\centering
\caption{The seven classes with examples and density at $N=10{,}000$.}\label{tab:classes}
\small
\setlength{\tabcolsep}{4pt}
\begin{tabular}{@{}llllr@{}}
\toprule
Class & Name & Definition & Examples & $N{=}10{,}000$ \\
\midrule
$S_0$ & Unit      & $n=1$               & 1       & 0.01\% \\
$S_1$ & Solar     & prime, $p\!\equiv\!1\,(6)$ & 7,13,19 & 6.1\% \\
$S_2$ & Lunar     & prime, $p\!\equiv\!5\,(6)$ & 5,11,17 & 6.2\% \\
$S_3$ & Dyadic    & $n=2$               & 2       & 0.01\% \\
$S_4$ & Triadic   & $n=3$               & 3       & 0.01\% \\
$S_5$ & Semiprime & $\Om(n)=2$          & 4,6,9   & 26.2\% \\
$S_6$ & Multiprime& $\Om(n)\geq 3$      & 8,12,24 & 61.5\% \\
\bottomrule
\end{tabular}
\end{table}

\begin{proposition}[Exhaustive partition]\label{prop:partition}
$\cl$ defines a partition of $\N$: every natural number belongs to
exactly one class.
\end{proposition}
\begin{proof}
The cases are hierarchically disjoint: $n=1,2,3$ tested first; remaining
primes classified by $\Om=1$ and residue mod~6 (Proposition~\ref{prop:prime-mod6}
guarantees all primes $>3$ land in $S_1$ or $S_2$); then
$\Om=2 \to S_5$, $\Om\geq 3 \to S_6$. \qed
\end{proof}

\begin{remark}[Density structure]\label{rem:density}
All classes except $S_6$ have natural density 0. For $N \geq 100$, more
than 60\% of integers are Multiprimes.
\end{remark}

\subsection{The Multiplicative Monoid}

\begin{definition}\label{def:otimes}
For $C_1, C_2 \in S$, define
$C_1 \otimes C_2 := \cl(n_1 \cdot n_2)$ for any $n_i \in C_i$.
\end{definition}

\begin{theorem}[Well-definedness]\label{thm:welldef}
The operation $\otimes$ is well-defined: $\cl(n_1 n_2)$ depends only on
$\cl(n_1)$ and $\cl(n_2)$.
\end{theorem}
\begin{proof}
The singleton classes $S_0=\{1\}$, $S_3=\{2\}$, $S_4=\{3\}$ have unique
representatives, so products involving them are trivially well-defined.
For non-singleton classes ($S_1, S_2, S_5, S_6$), note that
$\Om(n_1 n_2) = \Om(n_1) + \Om(n_2) \geq 2$ by Proposition~\ref{prop:omega-add}, so the
product never falls into the $\Om=1$ range where the mod-6 residue
distinction applies. The output class is therefore determined entirely by
the $\Om$-sum: $\Om = 2 \to S_5$; $\Om \geq 3 \to S_6$. Since $\Om$ is
the sole determinant for classes with $\Om \geq 2$, the map
$\cl(n_1 n_2)$ depends only on $\cl(n_1)$ and $\cl(n_2)$. \qed
\end{proof}

\begin{theorem}[Complete Cayley table]\label{thm:cayley}
The multiplication $\otimes$ on $S$ is given by Table~\ref{tab:cayley}.
\end{theorem}

%% Wide table → spans both columns
\begin{table*}[htbp]
\centering
\caption{Cayley table of $(S, \otimes)$. Verified on 200 pairs per
         entry (9{,}800 total), zero deviations (E7.1).}\label{tab:cayley}
\small
\setlength{\tabcolsep}{5pt}
\begin{tabular}{c|ccccccc}
\toprule
$\otimes$ & $S_0$ & $S_1$ & $S_2$ & $S_3$ & $S_4$ & $S_5$ & $S_6$ \\
\midrule
$S_0$ & $S_0$ & $S_1$ & $S_2$ & $S_3$ & $S_4$ & $S_5$ & $S_6$ \\
$S_1$ & $S_1$ & $S_5$ & $S_5$ & $S_5$ & $S_5$ & $S_6$ & $S_6$ \\
$S_2$ & $S_2$ & $S_5$ & $S_5$ & $S_5$ & $S_5$ & $S_6$ & $S_6$ \\
$S_3$ & $S_3$ & $S_5$ & $S_5$ & $S_5$ & $S_5$ & $S_6$ & $S_6$ \\
$S_4$ & $S_4$ & $S_5$ & $S_5$ & $S_5$ & $S_5$ & $S_6$ & $S_6$ \\
$S_5$ & $S_5$ & $S_6$ & $S_6$ & $S_6$ & $S_6$ & $S_6$ & $S_6$ \\
$S_6$ & $S_6$ & $S_6$ & $S_6$ & $S_6$ & $S_6$ & $S_6$ & $S_6$ \\
\bottomrule
\end{tabular}
\end{table*}

\begin{theorem}[Monoid structure {\cite{CliffordPreston1961}}]\label{thm:monoid}
$(S, \otimes)$ is a finite commutative monoid:
\begin{enumerate}[noitemsep, topsep=2pt]
\item \emph{Associativity:}
      $(C_1 \otimes C_2) \otimes C_3 = C_1 \otimes (C_2 \otimes C_3)$.
\item \emph{Neutral element:} $S_0 \otimes C = C$ for all $C$.
\item \emph{Absorbing element:} $S_6 \otimes C = S_6$ for all $C$.
\item \emph{Commutativity:} $C_1 \otimes C_2 = C_2 \otimes C_1$.
\item \emph{Idempotents:} $\{S_0, S_6\}$ only.
\end{enumerate}
\end{theorem}
\begin{proof}
Properties (2)--(5) are verified by inspection of
Table~\ref{tab:cayley}. For~(1): pick representatives $a \in C_1$,
$b \in C_2$, $c \in C_3$. Then
$(C_1 \otimes C_2) \otimes C_3 = \cl(r \cdot c)$ for any
$r \in \cl(a \cdot b)$. By Theorem~\ref{thm:welldef}, $\cl(r \cdot c)$
depends only on $\cl(r)$ and $\cl(c)$, so
$\cl(r \cdot c) = \cl((a \cdot b) \cdot c) = \cl(a \cdot (b \cdot c))$,
where the last equality uses associativity of integer multiplication.
By the symmetric argument, $C_1 \otimes (C_2 \otimes C_3) = \cl(a \cdot (b \cdot c))$. \qed
\end{proof}

\begin{theorem}[Completeness of the 7-class partition]
\label{thm:maximality}
The classification $\cl$ exhausts all distinctions available from the two
invariants $\Om(n)$ and $p \bmod 6$:
\begin{enumerate}[noitemsep, topsep=2pt]
\item \emph{Closure:} The partition is closed under the induced
      multiplication $\otimes$ (Theorem~\ref{thm:welldef}).
\item \emph{Completeness:} The decision tree branches on three levels:
      (i)~$\Om$-depth ($0, 1, 2, \geq 3$), (ii)~singleton status
      ($n=1,2,3$), and (iii)~residue mod~6 for primes $>3$. Since
      $\Om$-depth is already separated into all distinguishable strata
      and mod-6 residues admit exactly two prime classes ($1$ and $5$),
      no finer partition is possible within these invariants.
\item \emph{Necessity of each class:} Merging $S_1$ and $S_2$ destroys
      the deterministic addition law $S_1 + S_2 \to S_6$
      (Theorem~\ref{thm:solar-lunar}); merging $S_3$ or $S_4$ into a
      generic prime class removes the singleton property that forces the
      six trivially deterministic addition laws.
\end{enumerate}
\emph{Note:} In the multiplicative direction alone, the four prime classes
$S_1, S_2, S_3, S_4$ are equivalent (identical rows in
Table~\ref{tab:cayley}), so a 4-element quotient
$\{S_0, S_\mathrm{prime}, S_5, S_6\}$ suffices. The full 7-class
refinement is necessary only when addition laws are considered.
\end{theorem}
\begin{proof}
Property~(1) holds by Theorem~\ref{thm:welldef}.
Property~(2) follows from the exhaustive case analysis: the three
branching levels yield exactly $3 + 2 + 2 = 7$ leaf classes, with no
further splitting possible.
Property~(3): $p + q \equiv 0 \pmod{6}$ in
Theorem~\ref{thm:solar-lunar} requires $p \equiv 1$ and $q \equiv 5$ (or
vice versa); collapsing $S_1, S_2$ eliminates this distinction.
Collapsing $S_3$ or $S_4$ eliminates the singleton sums. \qed
\end{proof}

\subsection{Power Dynamics}

\begin{theorem}[Exponentiation trajectories]\label{thm:powerdyn}
For $n \in C$ and $k \geq 1$: $\Om(n^k) = k \cdot \Om(n)$.
\end{theorem}

The resulting trajectories are summarised in Table~\ref{tab:powerdyn}.

\begin{table}[htbp]
\centering
\caption{Power dynamics $\cl(n^k)$ by starting class.}\label{tab:powerdyn}
\small
\setlength{\tabcolsep}{4pt}
\begin{tabular}{@{}lccc@{}}
\toprule
Class $C$ & $k=1$ & $k=2$ & $k \geq 3$ \\
\midrule
$S_0$                       & $S_0$ & $S_0$ & $S_0$ \\
Primes ($S_1,S_2,S_3,S_4$) & $C$   & $S_5$ & $S_6$ \\
$S_5$                       & $S_5$ & $S_6$ & $S_6$ \\
$S_6$                       & $S_6$ & $S_6$ & $S_6$ \\
\bottomrule
\end{tabular}
\end{table}

\begin{theorem}[Universality of Multiprime]\label{thm:universality}
For every $n > 1$, there exists $k \in \N$ such that $\cl(n^k) = S_6$.
\end{theorem}
\begin{proof}
Set $k = \lceil 3/\Om(n) \rceil$. Then
$\Om(n^k) = k \cdot \Om(n) \geq 3$, giving $\cl(n^k) = S_6$. \qed
\end{proof}

\noindent\textbf{Verification.} 99{,}999 integers $n \in [2, 100{,}000]$:
all reach $S_6$; maximum observed $k = 3$, matching the theoretical
maximum $\lceil 3/1 \rceil = 3$ for primes.

\smallskip\noindent\textbf{Interpretation.} Exponentiation drives every
number toward the Multiprime basin. The flow
$S_0 \to \text{Primes} \to S_5 \to S_6$ is irreversible under repeated
self-multiplication.

%% ────────────────────────────────────────────────────────────────────────────
\section{Addition Laws}
%% ────────────────────────────────────────────────────────────────────────────

\subsection{Deterministic Laws}

\begin{definition}\label{def:det-add}
An addition law $\{C_1, C_2\} \to C_3$ is \emph{mathematically
deterministic} if $\cl(n_1 + n_2) = C_3$ for \textbf{all}
$n_1 \in C_1$, $n_2 \in C_2$ without exception.
\end{definition}

\begin{theorem}[The 7 deterministic addition laws]\label{thm:7laws}
Exactly 7 unordered class pairs yield a deterministic result (Table~\ref{tab:7laws}).
\end{theorem}

\begin{table}[htbp]
\centering
\caption{The 7 deterministic addition laws.}\label{tab:7laws}
\small
\setlength{\tabcolsep}{3pt}
\begin{tabular}{@{}lll p{2.8cm}@{}}
\toprule
$C_1$ & $C_2$ & Result & Justification \\
\midrule
$S_0$ & $S_0$ & $S_3$ & $1+1=2$ (singleton) \\
$S_0$ & $S_3$ & $S_4$ & $1+2=3$ (singleton) \\
$S_0$ & $S_4$ & $S_5$ & $1+3=4=2^2$ (singleton) \\
$S_3$ & $S_3$ & $S_5$ & $2+2=4=2^2$ (singleton) \\
$S_3$ & $S_4$ & $S_2$ & $2+3=5$ (Lunar, singleton) \\
$S_4$ & $S_4$ & $S_5$ & $3+3=6=2{\times}3$ (singleton) \\
$\mathbf{S_1}$ & $\mathbf{S_2}$ & $\mathbf{S_6}$ & \textbf{Thm.~\ref{thm:solar-lunar}} \\
\bottomrule
\end{tabular}
\end{table}

\emph{Note.} The first six laws are deterministic because
$S_0=\{1\}$, $S_3=\{2\}$, $S_4=\{3\}$ are singletons. The seventh is
the unique non-trivially deterministic law.

\emph{Count.} There are $\binom{7}{2}+7=28$ unordered class pairs.
Non-determinism of each of the remaining 21 is established by explicit
counterexample.\footnote{For instance, $\{S_0, S_2\}$ might appear
deterministic empirically ($P(S_6) \approx 62.5\%$ in Table~\ref{tab:prob-add}), but is
refuted by a single witness: $1 + 5 = 6$ with $\cl(6) = S_5$, while
$1 + 11 = 12$ with $\cl(12) = S_6$. Experiment E-S0.S2 in the
repository confirms that every non-deterministic pair admits at least
two distinct output classes.}
Table~\ref{tab:prob-add} gives the full empirical distributions.

\begin{theorem}[Solar + Lunar $\to$ Multiprime]
\label{thm:solar-lunar}
For any Solar prime $p \in S_1$ and Lunar prime $q \in S_2$:
$\cl(p + q) = S_6$.
\end{theorem}
\begin{proof}
$p \equiv 1 \pmod{6}$ and $q \equiv 5 \pmod{6}$ imply
$p + q \equiv 0 \pmod{6}$, so $6 \mid (p+q)$.
Since $p \geq 7$ and $q \geq 5$, we have $p+q \geq 12$.
Write $p+q = 6m$ with $m \geq 2$. Then
$\Om(p+q) = \Om(6m) = \Om(6)+\Om(m) = 2 + \Om(m) \geq 3$,
giving $\cl(p+q) = S_6$. \qed
\end{proof}

\noindent\textbf{Verification (E4.1).} All 90{,}000 Solar--Lunar pairs
in $[7, 100{,}000]$: zero exceptions (95\% CI for $P(S_6)$:
$[0.99996, 1.0]$).

\subsection{Probabilistic Laws}

For the remaining 21 class pairs, $\cl(n_1 + n_2)$ varies with
representative choice. Table~\ref{tab:prob-add} characterises all 21 distributions
empirically.

\begin{table*}[htbp]
\centering
\caption{Empirical distributions for all 21 probabilistic addition laws
         ($n=100{,}000$ pairs each, 95\% CI width ${<}0.007$).}\label{tab:prob-add}
\small
\setlength{\tabcolsep}{5pt}
\begin{tabular}{@{}llrrr@{}}
\toprule
$C_1$ & $C_2$ & $P(S_6)$ & $P(S_5)$ & $P(\text{prime})$ \\
\midrule
$S_0$ & $S_1$ & 62.2\% & 31.1\% & 6.7\% \\
$S_0$ & $S_2$ & 62.5\% & 31.0\% & 6.5\% \\
$S_0$ & $S_5$ & 74.4\% & 19.2\% & 6.4\% \\
$S_0$ & $S_6$ & 71.6\% & 20.8\% & 7.6\% \\
$S_1$ & $S_1$ & 88.4\% & 11.6\% & ${<}0.1$\% \\
$S_1$ & $S_3$ & 82.5\% & 17.5\% & 0\% \\
$S_1$ & $S_4$ & 86.8\% & 13.2\% & 0\% \\
$S_1$ & $S_5$ & 88.3\% & 10.8\% & 0.9\% \\
$S_1$ & $S_6$ & 86.5\% & 12.3\% & 1.2\% \\
$S_2$ & $S_2$ & 88.4\% & 11.6\% & ${<}0.1$\% \\
$S_2$ & $S_3$ & 33.8\% & 42.3\% & 23.9\% \\
$S_2$ & $S_4$ & 79.8\% & 18.9\% & 1.3\% \\
$S_2$ & $S_5$ & 87.4\% & 11.2\% & 1.4\% \\
$S_2$ & $S_6$ & 86.5\% & 12.2\% & 1.3\% \\
$S_3$ & $S_5$ & 74.1\% & 25.9\% & 0\% \\
$S_3$ & $S_6$ & 72.5\% & 21.0\% & 6.5\% \\
$S_4$ & $S_5$ & 79.2\% & 17.3\% & 3.5\% \\
$S_4$ & $S_6$ & 72.8\% & 20.7\% & 6.5\% \\
$S_5$ & $S_5$ & 87.5\% & 12.5\% & 0\% \\
$S_5$ & $S_6$ & 87.0\% & 11.5\% & 1.5\% \\
$S_6$ & $S_6$ & 71.8\% & 20.6\% & 7.6\% \\
\bottomrule
\end{tabular}
\end{table*}

The symmetric Solar/Lunar structure
($P(S_6 \mid S_1{+}S_1) \approx P(S_6 \mid S_2{+}S_2) \approx 88.4\%$)
confirms the Dirichlet balance (Section~2.2).
The anomalous distribution of $S_2 + S_3$ (only 33.8\% $S_6$, 23.9\%
prime) arises from the special role of 2: adding 2 to a Lunar prime
$q \equiv 5 \pmod{6}$ gives $q+2 \equiv 1 \pmod{6}$, frequently yielding
another prime.

%% ────────────────────────────────────────────────────────────────────────────
\section{SpectralAddress: Numbers as Coordinates}
%% ────────────────────────────────────────────────────────────────────────────

\subsection{Definition}

\begin{definition}[SpectralAddress]\label{def:sa}
Fix a prime basis $\B = (p_1, \ldots, p_K)$ with $p_i$ the $i$-th prime.
The SpectralAddress of $n$ is:
\[
\SA(n) = \bigl(\nu_{p_1}(n), \ldots, \nu_{p_K}(n)\bigr)
         \in \mathbb{N}_0^K
\]
Default basis: $K=7$, $\B=(2,3,5,7,11,13,17)$.
\end{definition}

\noindent\textbf{Limitation.} SpectralAddress is lossless only for
$\B$-smooth integers~\cite{Hildebrand1986}. For $K=7$ and $n$ up to
$17^7 = 410{,}338{,}673$, the representation is exact.

\subsection{Algebraic Isomorphisms}

\begin{proposition}[Multiplication isomorphism]\label{prop:mult-isom}
For $\B$-smooth $a, b$:
\[
\SA(ab) = \SA(a) + \SA(b)
\quad \text{(component-wise, $O(K)$)}
\]
\end{proposition}

\begin{proposition}[GCD and LCM isomorphisms]\label{prop:gcd-lcm-isom}
For $\B$-smooth $a, b$:
\begin{align*}
\SA(\gcd(a,b)) &= \min(\SA(a), \SA(b)) \quad \text{($O(K)$)} \\
\SA(\mathrm{lcm}(a,b)) &= \max(\SA(a), \SA(b)) \quad \text{($O(K)$)}
\end{align*}
\end{proposition}

\begin{proposition}[Coprimality criterion]\label{prop:coprime}
For $\B$-smooth $a, b$:
\[
\gcd(a,b) = 1 \iff \SA(a) \cdot \SA(b) = 0
\]
\textbf{Warning (scope restriction).} This criterion holds \emph{only}
for $\B$-smooth integers. If $a$ or $b$ carries a prime factor outside
$\B$, the SA-vectors cannot detect it, yielding a false positive.
\emph{Example:} $a = 2 \cdot 19 = 38$, $b = 3 \cdot 19 = 57$ with the
default basis give inner product $0$, yet $\gcd(38,57)=19\neq 1$.
\end{proposition}

\begin{table}[htbp]
\centering
\caption{Arithmetic operations as geometry.}
\small
\setlength{\tabcolsep}{4pt}
\begin{tabular}{@{}lll@{}}
\toprule
Arithmetic operation & SA operation & Complexity \\
\midrule
Multiplication $a \times b$  & Vector addition        & $O(K)$ \\
$\gcd(a, b)$                 & Component-wise min     & $O(K)$ \\
$\mathrm{lcm}(a, b)$         & Component-wise max     & $O(K)$ \\
Coprimality test             & Inner product $= 0?$   & $O(K)$ \\
Classification $\cl(n)$      & $L_1$ norm of $\SA(n)$ & $O(K)$ \\
\bottomrule
\end{tabular}
\end{table}

\noindent\textbf{Verification.} Bijection: 10{,}000/10{,}000 round-trip
reconstructions exact. Four isomorphisms: 40{,}000/40{,}000 tests passed.

%% ────────────────────────────────────────────────────────────────────────────
\section{Resonance Relation and Sparsity}
%% ────────────────────────────────────────────────────────────────────────────

\subsection{Definition}

\begin{definition}[Resonance]\label{def:resonance}
For integers $i, j$:
\[
R(i,j) = \|\min(\SA(i), \SA(j))\|_1 = \Om(\gcd(i,j))
\]
We call $i$ and $j$ \emph{resonant} if $R(i,j) > 0$, i.e.,
$\gcd(i,j) > 1$. The term \emph{resonance} denotes structural
co-divisibility; it has no topological or analytic content beyond this
arithmetic definition.
\end{definition}

\subsection{Asymptotic Sparsity}

\begin{theorem}[Coprime density]\label{thm:coprime-density}
\[
P\bigl(R(i,j) = 0\bigr)
= P\bigl(\gcd(i,j) = 1\bigr)
= \frac{6}{\pi^2} \approx 0.6079
\]
\end{theorem}
\begin{proof}
Classical:
$P(\gcd(i,j)=1) = \prod_p (1 - p^{-2}) = 1/\zeta(2) = 6/\pi^2$. \qed
\end{proof}

\noindent\textbf{Verification.} Empirical coprime fractions for
$T \in \{64, \ldots, 8192\}$ converge monotonically to $6/\pi^2$
(deviation at $T=8{,}192$: $0.0002$).

\subsection{Multi-Scale Resonance Hierarchy}

Experiment E6.2 ($T = 2{,}100$; 2{,}203{,}950 total pairs):

\begin{table}[htbp]
\centering
\caption{Multi-scale resonance hierarchy (E6.2).}\label{tab:resonance}
\small
\setlength{\tabcolsep}{4pt}
\begin{tabular}{@{}llll@{}}
\toprule
Type    & Definition      & Empirical         & Theoretical \\
\midrule
Type-2    & both div.\ 2   & 0.2499            & $\approx 1/4$     \\
Type-23   & both div.\ 6   & 0.02771           & $\approx 1/36$    \\
Type-235  & both div.\ 30  & 0.001096          & $\approx 1/900$   \\
Type-2357 & both div.\ 210 & $2.04\!\times\!10^{-5}$ & $\approx 1/44100$ \\
\bottomrule
\end{tabular}
\end{table}

This logarithmic hierarchy constitutes a natural multi-scale structure
arising from arithmetic. Experimental validation as a routing prior for
sequence processing is ongoing and will be presented in a companion paper.

%% ────────────────────────────────────────────────────────────────────────────
\section{Computational Verification}
%% ────────────────────────────────────────────────────────────────────────────

\subsection{Python Reference Implementation}

All theoretical results are verified in Python via a shared
classification module (\texttt{spectral\_matrix.py}); only E4.1 requires
\texttt{sympy} for symbolic verification. All 18 pytest-runnable
experiments are listed in the repository's \texttt{EXPERIMENT\_INDEX.md}.

\begin{itemize}[noitemsep, topsep=2pt]
\item E0.0 (SpectralMemory benchmark): 2-layer, 4-head architecture on
      TinyShakespeare; 5 tests (context scaling, latency, memory,
      infinite retrieval, sliding-window comparison); see Section~\ref{subsec:structural}
\item E1.0 (Maximality of 7-class partition): confirmed that all 7 classes
      are necessary and no finer $\Om$/mod-6 partition exists
\item E1.1 ($\Om$-additivity): 100{,}000 random pairs, rate 1.0
\item E1.2 (Prime residues mod~6): all 664{,}579 primes below $10^7$
\item E2.1 (Partition completeness): verified exhaustive coverage
\item E2.3a (Coprime sparsity): empirical $6/\pi^2$ convergence
\item E2.3b (Monoid identity/absorption): verified $S_0$, $S_6$ properties
\item E3.1 (Power dynamics): 99{,}999 integers, max $k=3$
\item E4.1 (Solar+Lunar, \texttt{sympy}): 90{,}000 pairs, zero exceptions
\item E5.1 (SA bijection): 10{,}000 round-trips, zero failures
\item E5.2 (SA isomorphisms): 40{,}000 tests, 100\% pass rate
\item E5.3a (Classification recovery from SA): 56{,}282 $\B$-smooth
      integers from 100{,}000 samples; $\cl(n)$ via direct factorization
      agrees with $\cl$ recovered from $\SA(n)$ in 100\% of cases
      (0 disagreements)
\item E5.3b ($\B$-smooth boundary): bijection breaks at $n=19$
      (first prime ${>}\max(\B)$); $\B$-smooth fraction decays from
      28.7\% at $N{=}1{,}000$ to 9.45\% at $N{=}10{,}000$ to 2.58\%
      at $N{=}100{,}000$
\item E6.2 (Resonance hierarchy): 2{,}203{,}950 pairs at $T=2{,}100$
\item E7.1 (Cayley table): 200 pairs/entry (9{,}800 total), zero deviations
\item E10.0 (All 28 addition laws): 100{,}000 pairs each, 7 deterministic
      confirmed, 21 non-deterministic with witnesses
\end{itemize}

\subsection{C++20 High-Performance Implementation}

\begin{table}[htbp]
\centering
\caption{C++20 benchmarks (Intel Core i5, \texttt{-O3}).}
\small
\setlength{\tabcolsep}{4pt}
\begin{tabular}{@{}lll@{}}
\toprule
Operation & Throughput & Complexity \\
\midrule
$\cl(n)$ classification & 0.84\,M ops/s   & $O(\sqrt{n})$ \\
Cayley lookup           & 475\,M ops/s$^{a}$ & $O(1)$ \\
SA computation          & 19.9\,M ops/s$^{b}$ & $O(K)$ \\
SA GCD                  & 10.7\,M ops/s$^{c}$ & $O(K)$ \\
SA round-trip           & 10{,}000/10{,}000 & --- \\
\bottomrule
\end{tabular}
\end{table}

\noindent{\small
$^{a}$ Precomputed class indices; 559.9\,M ops/s on AMD Ryzen~9~3900
(GCC~14.2, \texttt{-O3 -march=native}).\\
$^{b}$ Includes trial division over $K=7$ basis primes.\\
$^{c}$ Standalone component-wise min on precomputed vectors; combined
SA-compute+GCD limited by $^{b}$. GCD step alone exceeds 1{,}000\,M
ops/s on precomputed vectors (AMD Ryzen~9~3900, verified).
}

\smallskip
The Cayley table fits in 49~bytes; a single lookup requires one array
access.

\subsection{Verilog RTL Proof of Concept}

A Verilog module \texttt{spectral\_cell.v}, intended for synthesis, implements
SpectralAddress lookup for $n < 1024$ via a precomputed BRAM18
lookup table ($1024 \times 21 = 21{,}504$~bits $\approx 2.6$\,KB):

\begin{lstlisting}[language=Verilog]
module spectral_cell #(parameter K=7) (
    input  clk,
    input  [9:0]   n_in,
    output reg [K*3-1:0] sa_out
);
  reg [K*3-1:0] lut [0:1023];
  initial $readmemb("sa_lut.mem", lut);
  always @(posedge clk) sa_out <= lut[n_in];
endmodule
\end{lstlisting}

\noindent\textbf{Proof-of-concept results:} 14/14 test vectors passed
(Python behavioral emulation of the RTL design). This constitutes a
functional proof of concept, not exhaustive hardware verification.
RTL synthesis, place-and-route, timing simulation, and full-coverage
testing (all 1{,}024 inputs) remain pending. Target platform:
Xilinx Artix-7 (Arty~A7-35T); latency of 1 clock cycle at 100\,MHz is
the design target, not a verified timing result.

%% ────────────────────────────────────────────────────────────────────────────
\section{Computational Framework}
%% ────────────────────────────────────────────────────────────────────────────

\subsection{Operational Summary}

The individual components ($\Om$-function, mod-6 residues, $p$-adic
valuations) are classical. The contribution of this framework is their
combination into a single classification scheme with a complete algebraic
characterization, deterministic addition laws, and systematic
operationalization as a computational toolkit.

\begin{table}[htbp]
\centering
\small
\setlength{\tabcolsep}{4pt}
\begin{tabularx}{\columnwidth}{@{}lXX@{}}
\toprule
Concept & Classical view & PrimeSpec oper. \\
\midrule
What is $n$?    & A count           & Coord.\ $\SA(n)$ \\
Connects $i,j$? & Diff.\ $|i{-}j|$ & Resonance $R(i,j)$ \\
What is prime?  & Not divisible     & Atom: 1 axis \\
Multiply?       & Rpt.\ addition    & Vec.\ add.\ $O(K)$ \\
What is GCD?    & Largest divisor   & Min component \\
\bottomrule
\end{tabularx}
\caption{Classical view vs.\ PrimeSpec operationalization.}
\end{table}

\subsection{Informational Content}

The classification entropy at $N=10{,}000$ (computed from the empirical
densities in Table~\ref{tab:classes}):
\[
H(S) = -\sum_{i=0}^{6} P(S_i) \log_2 P(S_i)
      \approx 1.47 \text{ bits per integer}
\]
This lies well below the maximum $\log_2 7 \approx 2.807$ bits, reflecting
the strong dominance of $S_6$ (61.5\% of all integers at this range).

The SpectralAddress at $K=7$ carries $\approx 6.17$ bits per vector
(Shannon entropy of the coordinate distribution over the first $10^4$
integers).

\subsection{Structural Computation}\label{subsec:structural}

\textbf{For associative memory.} A query at position $q$ resonates with
all positions $p$ where $\gcd(q,p) > 1$. With a precomputed
SpectralAddress table, computing resonance scores for $N$ candidate
positions requires $O(N \cdot K)$ arithmetic operations.

\noindent\textbf{Experiment E0.0 (SpectralMemory benchmark).}
A 2-layer, 4-head SpectralMemory architecture replaces learned
positional encodings with $p$-adic valuation vectors and replaces the
causal attention mask with the resonance AND-gate
$R(i,j) = \sum_p \min(\nu_p(i), \nu_p(j))$. No positional parameters
are learned; all routing derives from prime structure. The architecture
is benchmarked against standard causal attention and a sliding-window
baseline on TinyShakespeare across five tests:

\begin{enumerate}[noitemsep, topsep=2pt]
\item \emph{Context scaling} ($T \in \{128, 256, 512\}$): The
      resonance-based mask achieves $\approx$61\% sparsity (matching
      the $6/\pi^2$ coprime fraction). The validation loss gap closes
      monotonically with context length, reaching parity at $T=512$
      ($\Delta_{\mathrm{val}} = {-}0.022$).
\item \emph{Latency scaling}: Wall-clock time per forward pass as a
      function of $T$.
\item \emph{Memory scaling}: Peak memory allocation as a function of $T$.
\item \emph{Infinite retrieval}: Resonance-based retrieval of a target
      token at distances $10{,}000$ to $500{,}000$ positions runs in
      $O(K)$ time independent of distance.
\item \emph{Sliding-window comparison}: Direct comparison with
      Longformer-style fixed-window attention.
\end{enumerate}

\noindent\emph{Status and scope.} These results demonstrate the
operational feasibility of resonance-based attention routing but are
limited to a small-scale setting (TinyShakespeare, 2-layer model, CPU).
Large-scale validation on production datasets and GPU hardware is
deferred to a companion paper. No performance superiority claims are
made in the present work.

\textbf{For database indexing.} Records indexed by $\B$-smooth key $n$
have SpectralAddress $\SA(n)$. Divisibility queries reduce to single-axis
lookups; structural similarity queries reduce to $K$-comparison resonance
tests. The $\B$-smoothness restriction must be verified for all keys.

\textbf{For computational hardware.} The entire Cayley table fits in
49~bytes; SpectralAddress lookup for $n < 1024$ fits in ${\approx}\,2.6$\,KB BRAM.
The resonance gate is a single bitwise AND. These operations are
hardware-native at commodity FPGA resource costs.

%% ────────────────────────────────────────────────────────────────────────────
\section{Related Work and Open Problems}\label{sec:relwork}
%% ────────────────────────────────────────────────────────────────────────────

\subsection{Positioning}

The individual components ($\Om$-function, mod-6 residues, $p$-adic
valuations, coprime density) are classical, and contributions C1, C2, C5
systematize known results into a unified framework.
The novel content is the complete characterization of deterministic
addition laws (C3), in particular the Solar+Lunar law
(Theorem~\ref{thm:solar-lunar}), and the combined operationalization as
a computational toolkit with cross-platform verification.
To our knowledge, no prior work has combined these components into a
single classification scheme with (1) a complete Cayley table, (2) a
full characterization of deterministic addition laws with explicit
non-determinism witnesses, or (3) cross-platform computational
verification including hardware behavioral simulation. Geroldinger and colleagues study
non-unique factorization in monoids~\cite{GeroldingerHalterKoch2006}.
$p$-adic applications in ML~\cite{ZunigaGalindo2023} use ultrametric
topology; our SpectralAddress provides a compatible finite-dimensional
substrate. For computational aspects of primality and factorization
see~\cite{CrandallPomerance2005}; for additive number theory
see~\cite{Nathanson1996}.

\subsection{Limitations}

\begin{enumerate}[noitemsep, topsep=2pt]
\item SpectralAddress is exact only for $\B$-smooth integers; non-smooth
      factors require extended basis or explicit treatment.
\item The monoid's essential structure is 4-element; the 7-class
      refinement is necessary for addition laws and SpectralAddress but
      not for multiplication alone.
\item ML application results are preliminary; large-scale validation is
      ongoing and deferred to a companion paper. No performance claims are
      made in the present work.
\end{enumerate}

\subsection{Open Problems}

\begin{enumerate}[noitemsep, topsep=2pt]
\item \textbf{Complete addition characterization.} Find closed-form
      distributions for the 21 non-deterministic addition laws.
\item \textbf{Finer partitions.} Does there exist a partition based on
      $\Om$ and residues mod $m > 6$ that is closed under multiplication
      and admits a strictly larger set of non-trivially deterministic
      addition laws?
\item \textbf{Gaussian integer extension.} Does an analogous semigroup
      exist over $\Z[i]$?
\item \textbf{Optimal basis size.} How does information content and
      computational utility scale with $K$?
\item \textbf{Prime pair sums in $S_6$.} Does every sufficiently large
      $n$ with $6 \mid n$ arise as a sum $p + q$ with $p \equiv 1\pmod{6}$
      and $q \equiv 5\pmod{6}$ both prime? This is \emph{strictly stronger}
      than Goldbach's conjecture restricted to multiples of~6, since it
      additionally imposes residue conditions on the summands. Standard
      circle-method results~\cite{Nathanson1996} apply to the unconstrained
      case; the residue-constrained variant appears open.
\item \textbf{$L$-function connections.} Are there relationships between
      the monoid structure and the analytic theory of the Selberg
      class~\cite{IwaniecKowalski2004}?
\end{enumerate}

%% ────────────────────────────────────────────────────────────────────────────
\section{Conclusion}
%% ────────────────────────────────────────────────────────────────────────────

We have completely characterised the algebraic structure induced by
classifying $\N$ via $\Om$-depth and prime residue mod~6.

\begin{table}[htbp]
\centering
\small
\setlength{\tabcolsep}{3pt}
\begin{tabularx}{\columnwidth}{@{}lXl@{}}
\toprule
Result & Statement & Status \\
\midrule
7-class partition & Exhaustive partition       & Proved \\
Monoid            & Finite comm.\ monoid       & Proved \\
Cayley table      & 49 entries verified         & Proved+ver. \\
Universality      & $n{>}1$ reaches $S_6$      & Proved+ver. \\
7 add.\ laws      & Exact, deterministic       & Proved \\
Solar+Lunar       & $S_1+S_2 \to S_6$          & Proved \\
SpectralAddress   & 4 isom., $O(K)$            & Proved+ver. \\
Resonance         & $6/\pi^2$ non-res.         & Proved+ver. \\
Verilog RTL       & PoC 14/14 vectors           & RTL pending \\
\bottomrule
\end{tabularx}
\caption{Summary of contributions and verification status.}
\end{table}

By making prime-arithmetic structure explicit, computable, and
hardware-native, we enable systems that reason about numbers not only as
ordered magnitudes but as elements of a structured algebraic space.

\smallskip
\noindent\textbf{Code and data:} \repourl

%% ────────────────────────────────────────────────────────────────────────────
\begin{thebibliography}{99}

\bibitem{Dirichlet1837}
P.G.L. Dirichlet,
``Beweis des Satzes, dass jede unbegrenzte arithmetische Progression
unendlich viele Primzahlen enthält,''
\textit{Abh. Königl. Preuss. Akad. Wiss.}, pp.\ 45--81, 1837.

\bibitem{ErdosKac1940}
P. Erd\H{o}s and M. Kac,
``The Gaussian law of errors in the theory of additive number theoretic
functions,''
\textit{Amer.\ J.\ Math.}, vol.\ 62, pp.\ 738--742, 1940.

\bibitem{GeroldingerHalterKoch2006}
A. Geroldinger and F. Halter-Koch,
\textit{Non-Unique Factorizations.} Chapman \& Hall/CRC, 2006.

\bibitem{Gouvea1997}
F.Q. Gouv\^{e}a,
\textit{p-adic Numbers: An Introduction.} Springer, 1997.

\bibitem{HardyWright1979}
G.H. Hardy and E.M. Wright,
\textit{An Introduction to the Theory of Numbers}, 5th~ed.
Oxford Univ.\ Press, 1979.

\bibitem{Sathe1953}
L. Sathe,
``On a problem of Hardy on the distribution of integers having a given
number of prime factors,''
\textit{J.\ Indian Math.\ Soc.}, vol.\ 17, pp.\ 63--141, 1953.

\bibitem{Selberg1954}
A. Selberg,
``Note on a paper by L.G.\ Sathe,''
\textit{J.\ Indian Math.\ Soc.}, vol.\ 18, pp.\ 83--87, 1954.

\bibitem{Tenenbaum2015}
G. Tenenbaum,
\textit{Introduction to Analytic and Probabilistic Number Theory},
3rd~ed. Cambridge Univ.\ Press, 2015.

\bibitem{ZunigaGalindo2023}
W.A. Z\'{u}\~{n}iga-Galindo,
``p-adic cellular neural networks,''
arXiv:2309.01512, 2023.

\bibitem{CliffordPreston1961}
A.H. Clifford and G.B. Preston,
\textit{The Algebraic Theory of Semigroups}, Vol.~1.
Amer.\ Math.\ Soc., 1961.

\bibitem{Nathanson1996}
M.B. Nathanson,
\textit{Additive Number Theory: The Classical Bases.}
Springer, 1996.

\bibitem{CrandallPomerance2005}
R. Crandall and C. Pomerance,
\textit{Prime Numbers: A Computational Perspective}, 2nd~ed.
Springer, 2005.

\bibitem{Hildebrand1986}
A. Hildebrand,
``On the number of positive integers $\leq x$ and free of prime
factors $> y$,''
\textit{J.\ Number Theory}, vol.\ 22, pp.\ 289--307, 1986.

\bibitem{IwaniecKowalski2004}
H. Iwaniec and E. Kowalski,
\textit{Analytic Number Theory.}
Amer.\ Math.\ Soc., 2004.

\end{thebibliography}

%% ────────────────────────────────────────────────────────────────────────────
\appendix
\section{Python Reference Classifier}
%% ────────────────────────────────────────────────────────────────────────────

\begin{lstlisting}[language=Python]
def omega(n):
    count, d = 0, 2
    while d * d <= n:
        while n % d == 0:
            count += 1; n //= d
        d += 1
    return count + (1 if n > 1 else 0)

def classify(n):
    if n == 1: return 'S0'
    if n == 2: return 'S3'
    if n == 3: return 'S4'
    w = omega(n)
    if w == 1: return 'S1' if n % 6 == 1 else 'S2'
    if w == 2: return 'S5'
    return 'S6'

def spectral_address(n, basis=(2,3,5,7,11,13,17)):
    sa = []
    for p in basis:
        v = 0
        while n % p == 0: v += 1; n //= p
        sa.append(v)
    return tuple(sa)

CAYLEY = [
    [0,1,2,3,4,5,6],
    [1,5,5,5,5,6,6],
    [2,5,5,5,5,6,6],
    [3,5,5,5,5,6,6],
    [4,5,5,5,5,6,6],
    [5,6,6,6,6,6,6],
    [6,6,6,6,6,6,6],
]
\end{lstlisting}

\end{document}
